% Figure: a phase diagram in the pressure-temperature plane, with the three
% coexistence curves meeting at the triple point and the liquid-vapour curve
% ending at the critical point. The slope of each curve is dp/dT = L/(T*dv).
\begin{figure}[htbp]
\centering
\begin{tikzpicture}
\begin{axis}[
  width=12cm, height=8cm,
  xlabel={temperature $T$}, ylabel={pressure $p$},
  xmin=0, xmax=10, ymin=0, ymax=10,
  axis lines=left,
  xtick=\empty, ytick=\empty,
  clip=false,
]
% sublimation curve (solid-vapour): from origin region up to triple point (4,2)
\addplot[engblue, very thick, smooth] coordinates {(1.2,0.15) (2.0,0.5) (3.0,1.1) (4.0,2.0)};
% melting curve (solid-liquid): steep, from triple point (4,2) up-left/up
\addplot[engred, very thick, smooth] coordinates {(4.0,2.0) (4.3,5.0) (4.6,9.5)};
% vaporisation curve (liquid-vapour): from triple point (4,2) to critical (8,7.5)
\addplot[engorange, very thick, smooth] coordinates {(4.0,2.0) (5.2,3.2) (6.4,4.9) (7.4,6.6) (8.0,7.5)};
% triple point
\filldraw[engblack] (axis cs:4.0,2.0) circle [radius=2pt];
\node[engblack, font=\scriptsize, anchor=north east] at (axis cs:4.0,2.0) {triple point};
% critical point
\filldraw[engblack] (axis cs:8.0,7.5) circle [radius=2pt];
\node[engblack, font=\scriptsize, anchor=south east] at (axis cs:8.0,7.5) {critical point};
% region labels
\node[enggray, font=\scriptsize] at (axis cs:2.2,6.5) {solid};
\node[enggray, font=\scriptsize] at (axis cs:6.2,8.7) {liquid};
\node[enggray, font=\scriptsize] at (axis cs:8.4,2.2) {vapour};
\end{axis}
\end{tikzpicture}
\caption[Coexistence curves and the Clausius-Clapeyron slope]{The three
coexistence curves of a pure substance on the pressure-temperature plane. Along
each curve two phases are in equilibrium, and the Clausius-Clapeyron relation
fixes the slope as $\mathrm{d}p/\mathrm{d}T = L/(T\,\Delta v)$, the latent heat
of the transition divided by the temperature and the volume change across it. The
liquid-vapour curve climbs steeply because the vapour is far more voluminous than
the liquid; the melting curve of most substances leans slightly right because the
liquid is slightly less dense than the solid, while water is the exception whose
melting curve leans left. The curve ends at the critical point, where the liquid
and vapour become indistinguishable and the latent heat falls to zero.}
\label{fig:vol04ch01:clausius-clapeyron}
\end{figure}
